% from ega3-1-fr.tex
\oldpage[III]{150}

\begin{lemma}[5.1.3.1]
\label{III.5.1.3.1-fr}
Sous les hypothèses de (\hyperref[III.5.1.1-fr]{5.1.1}), si on suppose en outre le
morphisme $f$ fermé, tout voisinage de $X'$ dans $X$ est identique à $X$.
\end{lemma}

\begin{proof}
Tout d'abord, on peut se ramener au cas où $f(X)=Y$. En effet, par hypothèse,
$f(X)$ est une partie fermée $Z$ de $Y$~; on peut en outre remplacer $f$ par
$f_{\mathrm{red}}$ (I, 6.3.4), et supposer par suite $X$ et $Y$ réduits~; on peut alors
remplacer $Y$ par le sous-préschéma fermé réduit de $Y$ ayant $Z$ pour espace
sous-jacent (I, 5.2.2), car tout idéal de $A$ est fermé, et tout anneau quotient de
$A$ est donc adique et noethérien. On a alors $f(X')=Y'$~; si $V$ est un voisinage
ouvert de $X'$ dans $X$, $f(X-V)$ est fermé dans $Y$ par hypothèse, et ne rencontre
pas $Y'$~; mais cela est impossible à moins que $X-V$ ne soit vide, puisque
$\mathfrak J$ est contenu dans le radical de $A$ (I, 1.1.15 et $0_{\mathrm I}$,
7.1.10), d'où la conclusion.
\end{proof}

Lorsqu'on se borne aux $\mathcal O_X$-Modules cohérents dont le support est propre
sur $Y$, (\hyperref[III.5.1.3-fr]{5.1.3}) peut s'énoncer, dans le langage des
catégories, en disant que le foncteur
$\mathcal F\mapsto\widehat{\mathcal F}$ est pleinement fidèle de la catégorie des
$\mathcal O_X$-Modules du type précédent, dans la catégorie des
$\mathcal O_{\widehat X}$-Modules cohérents, et établit par suite une équivalence de
la première de ces catégories avec une sous-catégorie pleine de la seconde
($0$, 8.1.6). Le théorème d'existence va prouver que lorsque $X$ est propre sur $Y$,
cette sous-catégorie est en fait la catégorie de tous les
$\mathcal O_{\widehat X}$-Modules cohérents. De façon précise~:

\begin{theorem}[5.1.4]
\label{III.5.1.4-fr}
Soient $A$ un anneau adique noethérien, $Y=\operatorname{Spec}(A)$,
$\mathfrak J$ un idéal de définition de $A$, $Y'=\mathrm V(\mathfrak J)$,
$f:X\to Y$ un morphisme séparé de type fini, $X'=f^{-1}(Y')$. Soient
$\widehat Y=Y_{/Y'}=\operatorname{Spf}(A)$, $\widehat X=X_{/X'}$ les complétés de
$Y$ et $X$ le long de $Y'$ et $X'$, $\widehat f:\widehat X\to\widehat Y$ le
prolongement de $f$ aux complétés~; alors, le foncteur
$\mathcal F\mapsto\mathcal F_{/X'}=\widehat{\mathcal F}$ est une équivalence de la
catégorie des $\mathcal O_X$-Modules cohérents de support propre sur
$\operatorname{Spec}(A)$, avec la catégorie des $\mathcal O_{\widehat X}$-Modules
cohérents de support propre sur $\operatorname{Spf}(A)$.
\end{theorem}

En d'autres termes, compte tenu de (\hyperref[III.5.1.3-fr]{5.1.3})~:

\begin{corollary}[5.1.5]
\label{III.5.1.5-fr}
Pour qu'un $\mathcal O_{\widehat X}$-Module soit isomorphe au complété d'un
$\mathcal O_X$-Module cohérent et de support propre sur $\operatorname{Spec}(A)$,
il faut et il suffit qu'il soit cohérent et de support propre sur
$\operatorname{Spf}(A)$.
\end{corollary}

Le cas le plus important est le

\begin{corollary}[5.1.6]
\label{III.5.1.6-fr}
Supposons $X$ propre sur $Y=\operatorname{Spec}(A)$. Alors le foncteur
$\mathcal F\mapsto\widehat{\mathcal F}$ est une équivalence de la catégorie des
$\mathcal O_X$-Modules cohérents et de la catégorie des
$\mathcal O_{\widehat X}$-Modules cohérents.
\end{corollary}

\begin{env}[5.1.7]
\label{III.5.1.7-fr}
\emph{Scholie.}
Si on tient compte de la caractérisation des faisceaux cohérents sur les préschémas
formels (I, 10.11.3), on voit que sous les conditions de
(\hyperref[III.5.1.1-fr]{5.1.1}), la donnée d'un $\mathcal O_X$-Module cohérent de
support propre sur $\operatorname{Spec}(A)$ équivaut (en posant
$Y_n=\operatorname{Spec}(A/\mathfrak J^{n+1})$ et $X_n=X\times_Y Y_n$) à la donnée
d'un système projectif de $\mathcal O_{X_n}$-Modules cohérents $(\mathcal F_n)$ tel
que pour $m\leq n$ on ait
\[
\mathcal F_m=\mathcal F_n\otimes_{\mathcal O_{Y_n}}\mathcal O_{Y_m}
\quad\text{(ou encore }\mathcal F_m=\mathcal F_n/\mathfrak J^{m+1}\mathcal F_n\text{)}
\]
et que le support de $\mathcal F_0$ soit une partie de $X_0$ propre sur $Y_0$. Au
moyen de (I, 10.11.4), on interprète de même les homomorphismes de
$\mathcal O_X$-Modules cohérents comme des homomorphismes de systèmes projectifs de
$\mathcal O_{X_n}$-Modules cohérents.

